% GENERATED FILE. Edit canonical lessons, then run npm run generate:books.
% canonical-source-hash: fnv1a64:2cd4714b53ba12fe
% canonical-lessons: JA-C445-zutsuu, JA-C445-fukutsuu, JA-C445-zukizuki, JA-C445-hakike, JA-C445-aza

\chapter{Aches and Pains}
\label{ch:ja-aches-and-pains}
\hlchaptermodality{445}

\begin{chapteropening}
\textbf{By the end of this chapter:} \emph{I can say where it hurts and how.}

Complete the last lesson of chapter 445: I can say where it hurts and how.
\end{chapteropening}

\section[\texorpdfstring{\ja{ずつう} (zutsū)}{ずつう (zutsū)}]{\ja{ずつう} (zutsū) — a headache}

\label{lesson:JA-C445-zutsuu}



\begin{quote}
Before the new one: say the Japanese for one's weight, then the Japanese for snoring.
\end{quote}

\subsection*{You'll want to know: \ja{ずつう}}
\textbf{\ja{ずつう}} — \emph{zutsū} — \textquotedblleft{}a headache\textquotedblright{}.

\begin{quote}
\textbf{\ja{けさから} \ja{ずつうが} \ja{します。}} — \emph{kesa kara zutsū ga shimasu} — \textquotedblleft{}I have had a headache since this morning.\textquotedblright{}
\end{quote}

\begin{cousinweb}
\textbf{\ja{ずつう}} (\emph{zutsū}) is Sino-Japanese (kango): two Chinese characters read with their on'yomi, readings borrowed from Chinese. They are \emph{zu}, \textquotedblleft{}head\textquotedblright{} (Middle Chinese \emph{duw}) and \emph{tsū}, \textquotedblleft{}pain\textquotedblright{} (\emph{thuwngH}) (KANJIDIC2; Middle Chinese in Baxter's transcription).
\end{cousinweb}

\begin{grammarlens}[title={What you've built}]
One more word for describing what you lived and felt.
\end{grammarlens}

\subsection*{Your turn}
\begin{itemize}
\raggedright
  \item \emph{Say it:} \emph{zutsū}
  \item \emph{Say it:} \emph{zutsū}, once more
  \item \emph{Say it:} \emph{kesa kara zutsū ga shimasu} — I have had a headache since this morning.
  \item \emph{Recall:} say the Japanese for one's weight, then the Japanese for snoring, then say \emph{zutsū} again
\end{itemize}

\begin{tcolorbox}[breakable,colback=teal!4,colframe=teal!35!black,title={Before you move on}]
What does \ja{ずつう} mean? (\textquotedblleft{}A headache\textquotedblright{}.) Say it once more.
\end{tcolorbox}
\section[\texorpdfstring{\ja{ふくつう} (fukutsū)}{ふくつう (fukutsū)}]{\ja{ふくつう} (fukutsū) — a stomach ache}

\label{lesson:JA-C445-fukutsuu}



\begin{quote}
Before the new one: say the Japanese for snoring, then the Japanese for a headache.
\end{quote}

\subsection*{You'll want to know: \ja{ふくつう}}
\textbf{\ja{ふくつう}} — \emph{fukutsū} — \textquotedblleft{}a stomach ache\textquotedblright{}.

\begin{quote}
\textbf{\ja{ふくつうで} \ja{がっこうを} \ja{やすみました。}} — \emph{fukutsū de gakkō o yasumimashita} — \textquotedblleft{}I stayed home from school with a stomach ache.\textquotedblright{}
\end{quote}

\begin{cousinweb}
\textbf{\ja{ふくつう}} (\emph{fukutsū}) is Sino-Japanese (kango): two Chinese characters read with their on'yomi, readings borrowed from Chinese. They are \emph{fuku}, \textquotedblleft{}abdomen\textquotedblright{} (Middle Chinese \emph{pjuwk}) and \emph{tsū}, \textquotedblleft{}pain\textquotedblright{} (\emph{thuwngH}) (KANJIDIC2; Middle Chinese in Baxter's transcription).
\end{cousinweb}

\begin{grammarlens}[title={What you've built}]
One more word for describing what you lived and felt.
\end{grammarlens}

\subsection*{Your turn}
\begin{itemize}
\raggedright
  \item \emph{Say it:} \emph{fukutsū}
  \item \emph{Say it:} \emph{fukutsū}, once more
  \item \emph{Say it:} \emph{fukutsū de gakkō o yasumimashita} — I stayed home from school with a stomach ache.
  \item \emph{Recall:} say the Japanese for snoring, then the Japanese for a headache, then say \emph{fukutsū} again
\end{itemize}

\begin{tcolorbox}[breakable,colback=teal!4,colframe=teal!35!black,title={Before you move on}]
What does \ja{ふくつう} mean? (\textquotedblleft{}A stomach ache\textquotedblright{}.) Say it once more.
\end{tcolorbox}
\section[\texorpdfstring{\ja{ずきずき} (zukizuki)}{ずきずき (zukizuki)}]{\ja{ずきずき} (zukizuki) — throbbing (of pain)}

\label{lesson:JA-C445-zukizuki}



\begin{quote}
Before the new one: say the Japanese for a headache, then the Japanese for a stomach ache.
\end{quote}

\subsection*{You'll want to know: \ja{ずきずき}}
\textbf{\ja{ずきずき}} — \emph{zukizuki} — \textquotedblleft{}throbbing (of pain)\textquotedblright{}.

\begin{quote}
\textbf{\ja{むしばが} \ja{ずきずき} \ja{いたみます。}} — \emph{mushiba ga zukizuki itamimasu} — \textquotedblleft{}My bad tooth is throbbing.\textquotedblright{}
\end{quote}

\begin{cousinweb}
\textbf{\ja{ずきずき}} (\emph{zukizuki}) is a native Japanese mimetic word, inherited rather than borrowed from Chinese, and its older origin is not settled (JMdict).
\end{cousinweb}

\begin{grammarlens}[title={What you've built}]
One more word for describing what you lived and felt.
\end{grammarlens}

\subsection*{Your turn}
\begin{itemize}
\raggedright
  \item \emph{Say it:} \emph{zukizuki}
  \item \emph{Say it:} \emph{zukizuki}, once more
  \item \emph{Say it:} \emph{mushiba ga zukizuki itamimasu} — My bad tooth is throbbing.
  \item \emph{Recall:} say the Japanese for a headache, then the Japanese for a stomach ache, then say \emph{zukizuki} again
\end{itemize}

\begin{tcolorbox}[breakable,colback=teal!4,colframe=teal!35!black,title={Before you move on}]
What does \ja{ずきずき} mean? (\textquotedblleft{}Throbbing (of pain)\textquotedblright{}.) Say it once more.
\end{tcolorbox}
\section[\texorpdfstring{\ja{はきけ} (hakike)}{はきけ (hakike)}]{\ja{はきけ} (hakike) — nausea; feeling sick}

\label{lesson:JA-C445-hakike}



\begin{quote}
Before the new one: say the Japanese for a stomach ache, then the Japanese for throbbing.
\end{quote}

\subsection*{You'll want to know: \ja{はきけ}}
\textbf{\ja{はきけ}} — \emph{hakike} — \textquotedblleft{}nausea; feeling sick\textquotedblright{}.

\begin{quote}
\textbf{\ja{ふねで} \ja{はきけが} \ja{しました。}} — \emph{fune de hakike ga shimashita} — \textquotedblleft{}I felt sick on the boat.\textquotedblright{}
\end{quote}

\begin{cousinweb}
\textbf{\ja{はきけ}} (\emph{hakike}) is a mixed word: native \emph{haki}, \textquotedblleft{}vomiting\textquotedblright{}, from \emph{haku}, joined to Sino-Japanese \emph{ke}, \textquotedblleft{}spirit, feeling\textquotedblright{} (Middle Chinese \emph{khj+jH}), a reading borrowed from Chinese (JMdict, KANJIDIC2; Middle Chinese in Baxter's transcription).
\end{cousinweb}

\begin{grammarlens}[title={What you've built}]
One more word for describing what you lived and felt.
\end{grammarlens}

\subsection*{Your turn}
\begin{itemize}
\raggedright
  \item \emph{Say it:} \emph{hakike}
  \item \emph{Say it:} \emph{hakike}, once more
  \item \emph{Say it:} \emph{fune de hakike ga shimashita} — I felt sick on the boat.
  \item \emph{Recall:} say the Japanese for a stomach ache, then the Japanese for throbbing, then say \emph{hakike} again
\end{itemize}

\begin{tcolorbox}[breakable,colback=teal!4,colframe=teal!35!black,title={Before you move on}]
What does \ja{はきけ} mean? (\textquotedblleft{}Nausea; feeling sick\textquotedblright{}.) Say it once more.
\end{tcolorbox}
\section[\texorpdfstring{\ja{あざ} (aza)}{あざ (aza)}]{\ja{あざ} (aza) — a bruise}

\label{lesson:JA-C445-aza}



\begin{quote}
Before the new one: say the Japanese for throbbing, then the Japanese for nausea.
\end{quote}

\subsection*{You'll want to know: \ja{あざ}}
\textbf{\ja{あざ}} — \emph{aza} — \textquotedblleft{}a bruise\textquotedblright{}.

\begin{quote}
\textbf{\ja{ひざに} \ja{あおい} \ja{あざが} \ja{あります。}} — \emph{hiza ni aoi aza ga arimasu} — \textquotedblleft{}There is a blue bruise on my knee.\textquotedblright{}
\end{quote}

\begin{cousinweb}
\textbf{\ja{あざ}} (\emph{aza}) is a native Japanese word (wago), inherited rather than borrowed from Chinese; its kanji spelling was attached to it later (JMdict).
\end{cousinweb}

\begin{grammarlens}[title={What you've built}]
One more word for describing what you lived and felt.
\end{grammarlens}

\subsection*{Your turn}
\begin{itemize}
\raggedright
  \item \emph{Say it:} \emph{aza}
  \item \emph{Say it:} \emph{aza}, once more
  \item \emph{Say it:} \emph{hiza ni aoi aza ga arimasu} — There is a blue bruise on my knee.
  \item \emph{Recall:} say the Japanese for throbbing, then the Japanese for nausea, then say \emph{aza} again
\end{itemize}

\begin{tcolorbox}[breakable,colback=teal!4,colframe=teal!35!black,title={Before you move on}]
What does \ja{あざ} mean? (\textquotedblleft{}A bruise\textquotedblright{}.) Say it once more.
\end{tcolorbox}
